\documentclass[12pt,a4paper,oneside]{article}
\usepackage[brazil]{babel}
\usepackage[utf8]{inputenc}
\usepackage{olymp}
\usepackage{color}
\usepackage[dvipsnames]{xcolor}
\usepackage{listings}

\setcounter{problem}{5}

\begin{document}

\begin{problem}{Imposto de Renda}{imposto}{2}

Em poucos dias termina o prazo para declaração anual do Imposto de Renda. O valor do imposto devido depende do total de vencimentos do ano passado. A seguinte tabela progressiva (a mesma há anos...) deve ser novamente usada para o cálculo do imposto a pagar.

\begin{table}[h!]
\centering
\begin{tabular}{ccc}
\hline
Base de cálculo anual em R\$ & Alíquota \% & Parcela a deduzir do imposto em R\$\\
\hline
Até 22.847,76 &- &-\\
De 22.847,77 até 33.919,80 &7,5 &1.713,58\\
De 33.919,81 até 45.012,60 &15,0 &4.257,57\\
De 45.012,61 até 55.976,16 &22,5 &7.633,51\\
Acima de 55.976,16 &27,5 &10.432,32\\
\hline
\end{tabular}
\end{table}

A tabela indica que quem recebeu até R\$ 22.847,76 não deve pagar imposto.

Quem recebeu acima disso e até R\$ 33.919,80 deve pagar 7,5\% do total recebido, deduzindo R\$1.713,58 deste valor. Assim, quem recebeu um total de
R\$30.000,00 deve pagar R\$ 641,01 de imposto, conforme cálculo a seguir:

\vspace{-6pt}
\begin{itemize}
\vspace{-6pt}
\item Aplicando a alíquota: 7,5\% de 30.000,00 = R\$ 2.250,00
\vspace{-6pt}
\item Deduzindo a parcela: R\$ 2.250,00 – R\$ 1.713,58 = R\$ 536,42
\end{itemize}
\vspace{-6pt}

O mesmo cálculo é feito nas demais faixas, aplicando a alíquota e deduzindo a parcela correspondente. Veja outro exemplo, para quem recebeu R\$ 50.000,00:

\vspace{-6pt}
\begin{itemize}
\vspace{-6pt}
\item Aplicando a alíquota: 22,5\% de 50.000,00 = R\$ 11.250,00
\vspace{-6pt}
\item Deduzindo a parcela: R\$ 11.250,00 – R\$ 7.633,51 = R\$ 3.616,49
\end{itemize}
\vspace{-6pt}

Faça um programa para calcular o imposto de renda devido. Seu programa deve obrigatoriamente usar uma função com o protótipo abaixo, que recebe o valor do total de vencimentos \texttt{v} e retorna o valor do imposto a pagar:

\texttt{\textcolor{Mulberry}{double} imposto(\textcolor{Mulberry}{double} v)}


%\Notes
%
%Observações, se tiver.

\InputFile

A entrada contém vários casos de teste. Cada caso de teste é dado em uma linha contendo um valor real $V$, que representa o total de vencimentos. Restrições: $0 \le V \leq 100.000.000$.

A entrada termina com $V = 0$, que não deve ser processada.

\OutputFile

A saída consiste em várias linhas, uma para cada caso de teste da entrada, contendo o valor do imposto a pagar, com duas casas decimais, seguindo o formato mostrado no exemplo a seguir.

% \Notes
% Nesta questão, seu programa deve usar a função \texttt{main} dada acima exatamente como está, não a modifique! Sua tarefa é apenas criar a função \texttt{eh\_primo}.

\Example

\begin{example}
\exmp{
30000
50000
12345
65432
30000.50
0
}{
R\$ 536.42
R\$ 3616.49
R\$ 0.00
R\$ 7561.48
R\$ 536.46
}%
\end{example}



%Comentários sobre o exemplo, se necessário.

\end{problem}

\end{document}
